% Options for packages loaded elsewhere
\PassOptionsToPackage{unicode}{hyperref}
\PassOptionsToPackage{hyphens}{url}
\PassOptionsToPackage{dvipsnames,svgnames,x11names}{xcolor}
%
\documentclass[
  10pt,
]{article}
\usepackage{amsmath,amssymb}
\usepackage{iftex}
\ifPDFTeX
  \usepackage[T1]{fontenc}
  \usepackage[utf8]{inputenc}
  \usepackage{textcomp} % provide euro and other symbols
\else % if luatex or xetex
  \usepackage{unicode-math} % this also loads fontspec
  \defaultfontfeatures{Scale=MatchLowercase}
  \defaultfontfeatures[\rmfamily]{Ligatures=TeX,Scale=1}
\fi
\usepackage{lmodern}
\ifPDFTeX\else
  % xetex/luatex font selection
  \setmainfont[]{DejaVu Serif}
  \setsansfont[]{DejaVu Sans}
  \setmonofont[]{DejaVu Sans Mono}
  \setmathfont[]{Latin Modern Math}
\fi
% Use upquote if available, for straight quotes in verbatim environments
\IfFileExists{upquote.sty}{\usepackage{upquote}}{}
\IfFileExists{microtype.sty}{% use microtype if available
  \usepackage[]{microtype}
  \UseMicrotypeSet[protrusion]{basicmath} % disable protrusion for tt fonts
}{}
\makeatletter
\@ifundefined{KOMAClassName}{% if non-KOMA class
  \IfFileExists{parskip.sty}{%
    \usepackage{parskip}
  }{% else
    \setlength{\parindent}{0pt}
    \setlength{\parskip}{6pt plus 2pt minus 1pt}}
}{% if KOMA class
  \KOMAoptions{parskip=half}}
\makeatother
\usepackage{xcolor}
\usepackage[margin=2.0cm]{geometry}
\usepackage{longtable,booktabs,array}
\usepackage{calc} % for calculating minipage widths
% Correct order of tables after \paragraph or \subparagraph
\usepackage{etoolbox}
\makeatletter
\patchcmd\longtable{\par}{\if@noskipsec\mbox{}\fi\par}{}{}
\makeatother
% Allow footnotes in longtable head/foot
\IfFileExists{footnotehyper.sty}{\usepackage{footnotehyper}}{\usepackage{footnote}}
\makesavenoteenv{longtable}
\setlength{\emergencystretch}{3em} % prevent overfull lines
\providecommand{\tightlist}{%
  \setlength{\itemsep}{0pt}\setlength{\parskip}{0pt}}
\setcounter{secnumdepth}{-\maxdimen} % remove section numbering
% doc_header.tex -- LaTeX preamble additions used by build_docs.sh for every
% Markdown -> LaTeX -> PDF conversion (supplements, README, submission files).
\usepackage{etoolbox}
% longtable may not be loaded by pandoc when a document has no table
\usepackage{longtable}
% bold vectors: \bm -> \symbf under unicode-math (XeLaTeX), bm package otherwise
\makeatletter
\@ifpackageloaded{unicode-math}{\providecommand{\bm}[1]{\symbf{#1}}}{\usepackage{bm}}
\makeatother
% smaller type in long tables so that wide audit tables fit the text width;
% inside a table a line break is allowed after an underscore of plain text
% (pandoc writes it as \_), so that long identifiers wrap inside their column
% rows are set a little apart, so that symbols with indices do not touch
\AtBeginEnvironment{longtable}{\footnotesize\renewcommand{\arraystretch}{1.45}%
  \let\sptunderscore\_\def\_{\sptunderscore\allowbreak}}
\setlength{\LTcapwidth}{\textwidth}
% allow line breaks after slashes and hyphens inside table cells
\usepackage{xurl}
% code blocks: break over-long lines instead of running into the margin.
% Plain ``` blocks are typeset by pandoc as {verbatim}, highlighted ones
% (```bash) as {Highlighting}; both are re-declared as breakable fancyvrb
% environments.  The continuation line is marked with a small hooked arrow.
\usepackage{fvextra}
\RecustomVerbatimEnvironment{verbatim}{Verbatim}{breaklines,breakanywhere,fontsize=\small}
\makeatletter
\@ifundefined{Highlighting}{}{%
  \RecustomVerbatimEnvironment{Highlighting}{Verbatim}{breaklines,breakanywhere,fontsize=\small,commandchars=\\\{\}}}
\makeatother
% inline code (pandoc writes it as \texttt{...}): never hyphenate -- an inserted
% hyphen could be read as part of a file name -- but allow a line break after
% an underscore, slash, comma, colon or equals sign, so that long paths and
% column lists wrap inside the text width and inside table columns.
% Inside a formula (e.g. the name of a notebook symbol in an equation) the same
% command must first leave math mode, which \text of amsmath does.
\ExplSyntaxOn
\tl_new:N \l__spt_code_tl
\cs_new_protected:Npn \__spt_texttt:n #1
  {
    \tl_set:Nn \l__spt_code_tl {#1}
    \regex_replace_all:nnN { ([\/\,\:\=]) } { \1 \c{allowbreak} } \l__spt_code_tl
    \regex_replace_all:nnN { \c{_} } { \c{_} \c{allowbreak} } \l__spt_code_tl
    \group_begin: \ttfamily \hyphenchar \font = -1 \scan_stop: \l__spt_code_tl \group_end:
  }
\RenewDocumentCommand \texttt { m }
  {
    \mode_if_math:TF { \text { \__spt_texttt:n {#1} } } { \__spt_texttt:n {#1} }
  }
\ExplSyntaxOff
% A figure is written in the Markdown source as an image paragraph followed by
% a caption paragraph ("**Figure S1.** ..."), so that the caption is also
% visible on GitHub.  An image that starts a paragraph is centered.
\usepackage{graphicx}% pandoc loads it only for documents that contain an image
\NewCommandCopy\sptincludegraphics\includegraphics
\RenewDocumentCommand\includegraphics{O{} m}{%
  \ifvmode\begin{center}\sptincludegraphics[#1]{#2}\end{center}%
  \else\sptincludegraphics[#1]{#2}\fi}
% keep figures where they are referenced
\usepackage{float}
\floatplacement{figure}{H}
% running header with the document title
\usepackage{fancyhdr}
\pagestyle{fancy}
\fancyhf{}
\fancyhead[L]{\small PROJECT spacetimes --- documentation and provenance}
\fancyhead[R]{\small\thepage}
\renewcommand{\headrulewidth}{0.4pt}
% 16x16 matrices: tighter columns so that a full matrix fits the text width
\setlength{\arraycolsep}{3.2pt}
% The documents of this project contain long tagged symbols, file paths and
% identifiers that cannot be hyphenated.  A paragraph is therefore allowed to
% have looser lines before TeX gives up and lets a line stick out of the margin.
\AtBeginDocument{\tolerance=2000 \emergencystretch=5em \hbadness=2000}
\ifLuaTeX
  \usepackage{selnolig}  % disable illegal ligatures
\fi
\IfFileExists{bookmark.sty}{\usepackage{bookmark}}{\usepackage{hyperref}}
\IfFileExists{xurl.sty}{\usepackage{xurl}}{} % add URL line breaks if available
\urlstyle{same}
\hypersetup{
  pdftitle={P2 --- The PDF files of the project folder: which were used, and how},
  pdfauthor={Prepared for Patrick Lee Nash, Ph.D.},
  colorlinks=true,
  linkcolor={blue},
  filecolor={Maroon},
  citecolor={blue},
  urlcolor={blue},
  pdfcreator={LaTeX via pandoc}}

\title{P2 --- The PDF files of the project folder: which were used, and
how}
\usepackage{etoolbox}
\makeatletter
\providecommand{\subtitle}[1]{% add subtitle to \maketitle
  \apptocmd{\@title}{\par {\large #1 \par}}{}{}
}
\makeatother
\subtitle{Provenance required by the task specification (PROJECT
spacetimes)}
\author{Prepared for Patrick Lee Nash, Ph.D.}
\date{November 12, 2026}

\begin{document}
\maketitle

\hypertarget{what-this-document-is-for}{%
\section{1. What this document is for}\label{what-this-document-is-for}}

The task specification asks for a provenance that lists every PDF file
that was used and says how it was used. The project folder contains nine
PDF files, all of them published papers. This document lists them, says
how each was read, what the project took from it, where the project
cites it, and what it did \emph{not} take. The symbols of these papers,
and the places where they collide with the symbols of the notebooks, are
the subject of the MASTER provenance
(\texttt{provenance/P1\_master\_symbols}, Sec. 5.2).

The result in one sentence: \textbf{no formula, no matrix and no number
of the project is copied from any of the nine papers}; the papers are
used for the origin and the conventions of the author's matrices, for
the mathematical background of the two spacetimes, and as references.
Everything that the project asserts about its matrices is proved by its
own programs (provenance ``dirac matrices'').

The nine files are not part of the repository, because they are
copyrighted publications. They are identified by their SHA-256 checksums
(Sec. 5), so that a reader who has the files can verify that he has the
same copies.

\hypertarget{the-nine-files}{%
\section{2. The nine files}\label{the-nine-files}}

``Pages'' is the number of pages of the file. ``Text'' says how the text
was obtained: \emph{digital} means that the file is the publisher's
digital original; \emph{layer} that it is a scan with a text layer made
by character recognition before the project (of mixed quality: words are
reliable, Greek letters and formulas often are not); \emph{OCR} that it
is a scan without text, which was read by optical character recognition
during this project (program tesseract 5.3.4) with the same limits.

\begin{longtable}[]{@{}
  >{\raggedright\arraybackslash}p{(\columnwidth - 8\tabcolsep) * \real{0.0460}}
  >{\raggedright\arraybackslash}p{(\columnwidth - 8\tabcolsep) * \real{0.3218}}
  >{\raggedright\arraybackslash}p{(\columnwidth - 8\tabcolsep) * \real{0.4828}}
  >{\raggedleft\arraybackslash}p{(\columnwidth - 8\tabcolsep) * \real{0.0690}}
  >{\raggedright\arraybackslash}p{(\columnwidth - 8\tabcolsep) * \real{0.0805}}@{}}
\toprule\noalign{}
\begin{minipage}[b]{\linewidth}\raggedright
Id
\end{minipage} & \begin{minipage}[b]{\linewidth}\raggedright
File
\end{minipage} & \begin{minipage}[b]{\linewidth}\raggedright
Paper
\end{minipage} & \begin{minipage}[b]{\linewidth}\raggedleft
Pages
\end{minipage} & \begin{minipage}[b]{\linewidth}\raggedright
Text
\end{minipage} \\
\midrule\noalign{}
\endhead
\bottomrule\noalign{}
\endlastfoot
D01 & \texttt{6dim\ dirac.pdf} & E. A. Lord, ``The Dirac spinor in six
dimensions'', Proc. Camb. Phil. Soc. 64, 765 (1968) & 14 & OCR \\
D02 & \texttt{A\ Remarkable\ Representation.pdf} & P. A. M. Dirac, ``A
remarkable representation of the 3 + 2 de Sitter group'', J. Math. Phys.
4, 901 (1963) & 10 & OCR \\
D03 & \texttt{A\ new\ spin-half.pdf} & P. L. Nash, ``A new spin-1/2 wave
equation'', J. Math. Phys. 26, 1363 (1985) & 5 & OCR \\
D04 & \texttt{Bryant-octon-17-2-185-txt.pdf} & R. L. Bryant,
``Submanifolds and special structures on the octonians'', J.
Differential Geom. 17, 185 (1982) & 48 & layer \\
D05 & \texttt{Calibrated\ geometries.pdf} & R. Harvey and H. B. Lawson,
Jr., ``Calibrated geometries'', Acta Math. 148, 47 (1982) & 111 &
layer \\
D06 & \texttt{Identities\ satisfied.pdf} & P. L. Nash, ``Identities
satisfied by the generators of the Dirac algebra'', J. Math. Phys. 25,
204 (1984) & 6 & OCR \\
D07 & \texttt{On\ the\ structure\ of\ the\ split\ octonion\ algebra.pdf}
& P. L. Nash, ``On the structure of the split octonion algebra'', Nuovo
Cimento B 105, 31 (1990) & 11 & layer \\
D08 & \texttt{Second\ gravity.pdf} & P. L. Nash, ``Second gravity'', J.
Math. Phys. 51, 042501 (2010) & 27 & digital \\
D09 & \texttt{Vector\ cross\ products.pdf} & R. B. Brown and A. Gray,
``Vector cross products'', Comment. Math. Helv. 42, 222 (1967) & 15 &
layer \\
\end{longtable}

Because of the limits of the recognized text, the project relies only on
statements of these papers that are legible \emph{in words} (for example
``the tilde denotes transpose'', or the ranges of the indices). Where a
formula could not be read with certainty, this is said (Sec. 3, D07).

\hypertarget{how-each-paper-was-used}{%
\section{3. How each paper was used}\label{how-each-paper-was-used}}

\textbf{D08, ``Second gravity'' (Nash 2010).} \emph{Read:} the passages
that define the space \(R^{4,4}\) and the matrices (its Sec. II), and
the list of references. \emph{Used for:} (i) the origin of the matrices
of \texttt{4+4.nb}: the paper defines real \(8\times8\) matrices
\(\tau^A\), \(\bar\tau^A\) with
\(\tau^A\bar\tau^B+\tau^B\bar\tau^A=2G_{AB}\) and the \(16\times16\)
matrices \(t^A\) with the blocks \((0,\bar\tau^A;\tau^A,0)\), which is
the construction of the notebook; (ii) the conventions that differ from
the notebook: the index runs from 1 to 8 with \(\tau^8=\bar\tau^8=1\),
the metric has the diagonal \((1,1,1,-1,-1,-1,-1,1)\), a tilde denotes
the transpose; (iii) the statement that the construction follows a
procedure of Lord (D01). \emph{Not used:} its matrices themselves (the
project takes the matrices from the notebook and did not compare them
entry by entry with the paper), and its physics (a theory of gravity).
\emph{Cited:} both manuscripts, Secs. I and II; P1, Secs. 5.2 and 6.

\textbf{D07, ``On the structure of the split octonion algebra'' (Nash
1990).} \emph{Read:} the passages that define the matrices (its Sec. 2).
\emph{Used for:} the same definitions as D08, in the same conventions;
its remark that, relative to an earlier paper of the author, ``tau bar''
and ``tau'' are exchanged and the sign of \(\tau^7\) is changed, which
is a warning for anyone who compares the papers with one another; its
condition that a product of the seven matrices with \(A=1,\dots,7\) be
\(+1\). In the scanned copy the formula of that product is not legible
enough to say which factors carry a bar; the project therefore states
its own result, \(\tau_{1644}^1\cdots\tau_{1644}^7=1\) (check
\texttt{44-B1}), and does not claim that this is the convention of the
paper. \emph{Not used:} the relation to the split octonion algebra and
to triality; the manuscripts refer to the paper for it and say so (``not
established here''). \emph{Cited:} both manuscripts, Secs. I and II; P1.

\textbf{D06, ``Identities satisfied by the generators of the Dirac
algebra'' (Nash 1984).} \emph{Read:} the passages that define its
matrices (its Sec. 2). \emph{Used for:} the origin of the six
antisymmetric \(4\times4\) matrices \(s_h\), \(t_h\) from which the
construction starts: the paper calls the six antisymmetric ones among
its real \(4\times4\) matrices \(s^h\) (self-dual) and \(t^h\)
(anti-self-dual), and the notebooks keep these letters
(\texttt{s4by4{[}h{]}}, \texttt{t4by4{[}h{]}}). \emph{Not used:} its
identities, which concern four-dimensional matrices. \emph{Cited:}
manuscripts, Sec. I and (4+4) Sec. II A; P1.

\textbf{D01, ``The Dirac spinor in six dimensions'' (Lord 1968).}
\emph{Read:} its introduction and the passages on the generators.
\emph{Used for:} one statement only: D08 says that its construction
follows the procedure of this paper, which builds spinor representations
of orthogonal groups from the generators of a Dirac algebra of lower
dimension. \emph{Cited:} manuscripts, Sec. I.

\textbf{D02, ``A remarkable representation of the 3 + 2 de Sitter
group'' (Dirac 1963).} \emph{Read:} introduction. \emph{Used for:} the
lineage of the real representations. The notebook
\texttt{Dirac\_Einstein-Lovelock\_4+4\_Nash.nb} says that the
calculation ``has its roots'' in this paper, in which Dirac defines a
real representation of the covering group of \(SO(3,2)\), and that the
author extended it to \(SO(3,3)\), to \(SO(4,4)\) and to real
\(16\times16\) matrices. The Introduction of the manuscripts repeats
this lineage in one sentence. \emph{Cited:} manuscripts, Sec. I.

\textbf{D03, ``A new spin-1/2 wave equation'' (Nash 1985).} \emph{Read:}
abstract and introduction. \emph{Used for:} the same lineage (the stage
of \(SO(3,3)\)), together with D06. \emph{Not used:} its wave equation.
\emph{Cited:} manuscripts, Sec. I.

\textbf{D04 (Bryant 1982), D05 (Harvey and Lawson 1982), D09 (Brown and
Gray 1967).} \emph{Read:} their introductions. \emph{Used for:} the
statement that the signature 7+1 is the case of the octonions, and as
references for the algebra of the octonions, the groups \(G_2\) and
\(Spin(7)\) and the vector cross products of seven and eight dimensions,
to which the 7+1 matrices are related. \emph{Not used:} any formula; the
octonionic matrices of the project come from \texttt{7+1.nb} and are
verified by machine (checks \texttt{71-B1}, \texttt{71-D6}). The
manuscripts say that the relation of the 7+1 matrices to this
mathematics is not established in them. \emph{Cited:} manuscripts, Secs.
I and II.

\hypertarget{statements-of-the-project-that-rest-on-these-papers}{%
\section{4. Statements of the project that rest on these
papers}\label{statements-of-the-project-that-rest-on-these-papers}}

So that a reader can check them against the papers, these are all the
statements of the project documents whose only support is one of the
nine files:

\begin{longtable}[]{@{}
  >{\raggedright\arraybackslash}p{(\columnwidth - 4\tabcolsep) * \real{0.6667}}
  >{\raggedright\arraybackslash}p{(\columnwidth - 4\tabcolsep) * \real{0.2361}}
  >{\raggedright\arraybackslash}p{(\columnwidth - 4\tabcolsep) * \real{0.0972}}@{}}
\toprule\noalign{}
\begin{minipage}[b]{\linewidth}\raggedright
Statement
\end{minipage} & \begin{minipage}[b]{\linewidth}\raggedright
Where
\end{minipage} & \begin{minipage}[b]{\linewidth}\raggedright
Paper
\end{minipage} \\
\midrule\noalign{}
\endhead
\bottomrule\noalign{}
\endlastfoot
The author's papers define matrices \(\tau^A\), \(\bar\tau^A\), \(t^A\)
of the kind used here, with the index values 1 to 8, the value 8
corresponding to the value 0 here & M II A; P1 5.2, 6.2 & D07, D08 \\
A tilde denotes the transpose in those papers & P1 5.2 & D07, D08 \\
D07 exchanges ``tau'' and ``tau bar'' and changes the sign of \(\tau^7\)
relative to an earlier paper & P1 5.2 & D07 \\
The construction follows a procedure of Lord & M I & D08, D01 \\
The letters \(s\), \(t\) of the \(4\times4\) matrices, self-dual and
anti-self-dual, are those of the 1984 paper & M II A (4+4); P1 5.2 &
D06 \\
The real representations descend from Dirac's representation of the
covering group of \(SO(3,2)\), through \(SO(3,3)\) & M I & D02, D03, D06
(and the text of the notebook N04) \\
The signature 7+1 is the case of the octonions & M I, II & D04, D05,
D09 \\
\end{longtable}

M denotes the two manuscripts (the section numbers are the same in both
unless a spacetime is named).

The manuscripts cite about seventy further works. These were not
supplied as files; they are cited from the published record, and their
bibliographic data are in \texttt{manuscript/spacetimes\_refs.bib}.

\hypertarget{checksums}{%
\section{5. Checksums}\label{checksums}}

The following lines of \texttt{provenance/inputs\_sha256.txt} identify
the nine copies (SHA-256; verify with
\texttt{shasum\ -a\ 256\ \textless{}file\textgreater{}} on macOS):

\begin{verbatim}
8d2a2c77c7d0fa0a864e9254d465771b6164355b863f969e561eb0677cfcdc90  6dim dirac.pdf
56e95fadec54e9b72e63f9beb0ddf6e424796622a64e7039a21b7ec98fc8fb3f  A Remarkable Representation.pdf
8eddd3a6a63d38531d881e12be56e350cd3f14846c112b53c59f7099ee6a86f1  A new spin-half.pdf
85a8fe5049e60b5b8a65bc531110bc371f96749360c120716b00c0daaebcc0cf  Bryant-octon-17-2-185-txt.pdf
af0befdd8949c78ca55f3c4aac8fb4f741083df9f23415cbd48c875160df64e3  Calibrated geometries.pdf
197d18bfe18917375b11703bbf0b2fe022d97afb41278024841cd0d0ead64d22  Identities satisfied.pdf
27fbbfb90522fecb4f23ec77a23ea621a87b0c9a3ad5d8b29f97dff2cd18d0bc  On the structure of the split octonion algebra.pdf
7be828aad7e05b1753f9b81fdb57f65983a8b38536b49f3bdafeaa16242fc36a  Second gravity.pdf
a436b29fd6b0e0d8ee9819398389619f4427b5384207bd404565ea569dde36c6  Vector cross products.pdf
\end{verbatim}

\end{document}
